\documentclass[12pt]{article}

%\usepackage[utf8]{inputenc}
%\usepackage[T2A]{fontenc}
%\usepackage[russian,english]{babel}
\usepackage{amsmath,amsfonts,amssymb,bm}

\setlength{\textwidth}{168mm}%210
\setlength{\textheight}{210mm}%297
\topmargin -10mm
\oddsidemargin  -1mm
\evensidemargin -1mm
\usepackage[utf8]{inputenc}
\usepackage{amsmath,amssymb}
\usepackage{latexsym}
\usepackage{slashed}
\usepackage{authblk}
\usepackage{xcolor}
\usepackage{hyperref}

\usepackage{cancel}
\usepackage{indentfirst}
\usepackage{graphicx}% Include figure filesF
\usepackage{epstopdf}
%\graphicspath{{figs/}}
\usepackage{dcolumn}% Align table columns on decimal point
\usepackage{bm}% bold math
\usepackage{lipsum}
\usepackage{slashed}
\usepackage{enumitem}
\usepackage{txfonts}
\usepackage{esint}
\usepackage{xcolor}
\usepackage{mathtools}
\usepackage{transparent}
\usepackage{upgreek}
\usepackage{tabularx}
\usepackage{csquotes}
\usepackage{siunitx}
%\usepackage{geometry}
%\geometry{a4paper,total={178mm,257mm},right=15mm,left=15mm,bottom=20mm,top=20mm}
%\usepackage[font=footnotesize,labelfont=bf,labelsep=period]{caption}
% \usepackage[bookmarks=true,colorlinks=true, allcolors=black,linkcolor=black]{hyperref}
%\usepackage{tikz}
%\usepackage[compat=1.1.0]{tikz-feynman}
%\tikzfeynmanset{/tikzfeynman/warn luatex=false}
%\usetikzlibrary{arrows,automata}
%\usetikzlibrary{matrix}
%\addto\captionsenglish{\renewcommand{\figurename}{Fig.}}

%\pagestyle{plain}%номера страниц
%\usepackage{setspace}%Полуторный
%\onehalfspacing% межстрочный интервал
%\doublespacing

%\usepackage{wrapfig}
% \usepackage[backend=biber,style=ieee,articletitle=false]{biblatex}
%%\usepackage[backend=biber, style=phys,articletitle=false,biblabel=brackets,
%%            chaptertitle=false,pageranges=false]{biblatex}
%%\ExecuteBibliographyOptions[misc]{eprint=true}
%%\ExecuteBibliographyOptions[online]{eprint=true}
%%\addbibresource{main.bib}

%\usepackage{definitions}

\begin{document}
\title{Hawking Radiation and Primordial Black Holes}

\author{
	Group № 2: Qichao Chang, Shuhua Hao, Hengbin Shao, Borodin Nikita, Konstantinova Liza, Artem Popov, Slautin Mikhail, Sekretov Mikhail, Kulikov Aleksey, Moskalenko Anastasia, Vikulova Anastasia, Kanzyuba Daria\\ Supervisor: Sergey Zavyalov \vskip 10mm \textit{XXV JINR-ISU Baikal Summer School on Physics of Elementary Particles and Astrophysics}}



\maketitle
\newpage
\section{Introduction}

\section{Conclusion}
\indent In this work, the various masses of the primordial black holes and their lifetimes were calculated. The work also investigated what would happen if there were a  black hole in the center of the Earth. It turned out that it would not be able to absorb all the matter around it. The maximum mass of a black hole that can have and still evaporate in the presence of the cosmic microwave background was calculated. The photon and neutrino energies of the primary black holes were calculated and compared with the range of measured energies of existing experimental facilities. It turned out that we need a new generation of installations with better resolution to study photons and neutrinos.


%\LaTeX{} \cite{MacGibbon_1998}
%\bibliographystyle{unsrt}
%\bibliography{ref}

\LaTeX{} \cite{MacGibbon_1998} is a set of macros built atop \TeX{} \cite{MacGibbon_1998}.
\bibliographystyle{plain} % We choose the "plain" reference style
\bibliography{ref} % Entries are in the refs.bib file

\end{document}
